\section*{Chapter \ref{ch:b3:institutions} --- Institutions and Long-Run Prosperity}

\begin{solution}{exo:b3:institutions:1}
$\ln(1{,}744/799)/28 = 2.8~\%$ and $\ln(19{,}238/1{,}744)/44 = 5.5~\%$ a year: about twice as fast. % figdata: opening_stats
\end{solution}

\begin{solution}{exo:b3:institutions:2}
$41{,}321/1{,}569 = 26.3$, or $\log_2 26.3 = 4.7$ doublings.
\end{solution}

\begin{solution}{exo:b3:institutions:3}
$-0.57/-0.61 = 0.93$: a unit more on the index raises log income by 0.93.
\end{solution}

\begin{solution}{exo:b3:institutions:4}
A gap of the same size before 1945, or one that opened only with differences in aid,
resources or war that the rules do not explain; or a North that grew as fast as the
South in periods when its rules did not change while the South's did.
\end{solution}

\begin{solution}{exo:b3:institutions:5}
$k^* = (a(1 - p)A)^2 = (1 - p)^2$: 1 at $p = 0$ and 0.49 at $p = 0.3$. Output $2\sqrt{k^*} = 2(1 - p)$: 2 and 1.4.
\end{solution}

\begin{solution}{exo:b3:institutions:6}
$-0.57/-0.20 = 2.85$, three times larger; the denominator is now small relative to its
standard error, so the estimate is imprecise, and if the \omterm{def:b3:instrumental-variables:tsls}{first stage} cannot be told
from zero the confidence interval is unbounded.
\end{solution}

\begin{solution}{exo:b3:institutions:7}
The tested solution is \code{ex07_transition.py}, in the chapter's \code{exercises/}
folder. Poland fell 15.8~\% to 1991 and was back in 1995; Czechia 9.3~\% to 1991, back in 1995; Russia 33.6~\% to
1996, back in 2003; Ukraine 54.3~\% to 1997, back in 2007. % output: ex07_transition
\end{solution}

\begin{solution}{exo:b3:institutions:8}
$(0.03 - 0.09) - (0.085 - 0.09) = -0.055$ thefts per block and month. Displacement would show as a rise on the
blocks next to the protected ones, above the trend of distant blocks.
\end{solution}

\begin{solution}{exo:b3:institutions:9}
A jump at the boundary in characteristics fixed before the mita --- elevation, slope,
ethnicity, pre-colonial settlement --- or people and firms sorting across the
boundary since; either would mean that the two sides differ by more than the mita.
\end{solution}

\begin{solution}{exo:b3:institutions:10}
$\operatorname{cov}(z, \ln y)/\operatorname{cov}(z, s) = \beta + \gamma\pi_h/\pi_s$. The bias is $\gamma\pi_h/\pi_s$: with high mortality going with worse \omterm{def:b3:institutions:institutions}{institutions} and
worse health ($\pi_s, \pi_h < 0$) and health raising income ($\gamma > 0$), it is positive.
\end{solution}

\begin{solution}{exo:b3:institutions:11}
The manager maximises $eR + (1 - e)bB - e^2/2$, so $e^* = R - bB$, which falls with $b$. With $R = 0.8$ and $B = 0.6$: 0.8 with no
bail-out, 0.5 at $b = 0.5$, 0.2 at $b = 1$.
\end{solution}

\begin{solution}{exo:b3:institutions:12}
Ten per cent more police change crime by $10 \times \varepsilon$ per cent, so 100 officers prevent $5{,}000\,|\varepsilon|$ crimes,
$50\,|\varepsilon|$ each. At $\varepsilon = 0$ they prevent none, and hiring never pays. At $-0.18$ each prevents 9 crimes, and
hiring pays if a crime costs society more than $80{,}000/9 = 8{,}889$ dollars; at $-0.5$ each prevents 25,
and the threshold is $80{,}000/25 = 3{,}200$ dollars. The verdict turns on the elasticity, which is
the contested number, and on the mix of crimes prevented: the costs per crime used in
such studies are far higher for violent crimes than for thefts, and the estimates that
hold up find larger effects on violent crime. % ledger: b3in:levitt, b3in:cm
\end{solution}

\begin{solution}{pb:b3:institutions:river}
% figdata: river_stats
\textbf{1.} North: 20{,}342, 25{,}120 and 61{,}423 dalens; South: 20{,}230, 25{,}281 and 40{,}420.

\textbf{2.} $-0.6$ log points in 1995 and 41.8 in 2025.

\textbf{3.} North 1.41~\% a year before and 2.98~\% after; South 1.49~\% and 1.56~\%.

\textbf{4.} 42.5 log points.

\textbf{5.} Parallel paths before 1995: 1.41 against 1.49~\% a year with a gap near zero --- close.

\textbf{6.} Anything else that changed in North in 1995 is counted with the titles; and
South may have been affected by North's reform, through trade or migration.

\textbf{7.} D2 0.393, D3 0.262, D4 0.032, D7 0.313; the others 0.

\textbf{8.} 0.34 log points.

\textbf{9.} 43.0 log points.

\textbf{10.} Close to the 42.5 of the difference in differences: both comparisons agree.

\textbf{11.} It is the neighbour most exposed to spillovers from North's reform, so it
may not be an untreated comparison.

\textbf{12.} That no combination of donors reproduces North's path, so the \omterm{def:b3:difference-in-differences:synthetic-control}{synthetic
control} is not credible after 1995 either.

\textbf{13.} 1.28.

\textbf{14.} $-0.00463$ per kilometre, $F = 37.4$.

\textbf{15.} $-0.00459$ per kilometre.

\textbf{16.} $-0.00459/-0.00463 = 0.99$.

\textbf{17.} Better soil raises both the titled share and income, an omitted variable
positively correlated with $s$.

\textbf{18.} The survey office was in a town: distance to it is also distance to
markets, schools and roads.

\textbf{19.} $\Delta s = 0.341$ and $\Delta\ln y = 0.438$: by 2SLS, $0.99 \times 0.341/0.438 = 77~\%$; by least squares, $1.28 \times 0.341/0.438 = 100~\%$.

\textbf{20.} $y_N/y_S = ((1 - p_N)/(1 - p_S))^{1/2}$.

\textbf{21.} The log gap is $0.99 \times 0.341 = 0.338$; $1 - p_S = 0.95e^{-2 \times 0.338}$, so $p_S = 0.52$.

\textbf{22.} A risk of losing half of all output is high; titles also make land
saleable, save the cost of guarding it and settle disputes.

\textbf{23.} Titled land can be pledged for loans, so titles raise investment through
credit as well; with that channel a smaller difference in risk would explain the gap,
and the pure expropriation model overstates the risk needed.

\textbf{24.} Property security, measured by the titled share and estimated by 2SLS, explains about 77~\% of the
44-log-point income gap between the banks (0.44 in natural-log units); least squares would attribute all of it.
\end{solution}
